% Einleger fuer den Spielleiterschirm, DIN A4 quer.
%
% Der Regellauf fuer dsa5einleger.cls: jedes Element genau einmal, und die
% drei Seiten nachgebaut, an denen die Klasse vermessen wurde.
%
%   cd beispiel
%   TEXINPUTS="..;" xelatex einleger.tex
%
% Zum Pruefen der Spaltenlage:
%   \documentclass[spaltenzeigen]{dsa5einleger}

\documentclass{dsa5einleger}

% Hilfslaenge fuer das letzte Beispiel.
\newlength{\dsahilfe}

\begin{document}

%%% ---------------------------------------------------------------- 0
%%% Eine Stimmungsbildseite: nur ein randfuellendes Bild mit umlaufendem
%%% Zierrahmen, so wie die ersten drei Seiten des offiziellen Einlegers.
%%% Das Bild wird deckend zugeschnitten, nicht gestreckt.
%%%
%%% Der Baukasten hat keine Illustration im Querformat, deshalb der
%%% Ruecksprung: liegt ein eigenes Bild in grafiken/, wird es genommen,
%%% sonst die neutrale Buchrueckseite. So baut der Regellauf bei jedem
%%% durch, der nur die Baukastengrafiken hat.
%%%
%%% Das Ideal ist ein Bild im Verhaeltnis 303:216 = 1,40. stimmung-dorf hat
%%% 2528 x 1696 px, also 1,49 — \dsaBildDeckend skaliert es auf die
%%% Blatthoehe und beschneidet die Breite mittig, je 9,5 mm links und
%%% rechts. Auf 303 mm Blattbreite steht es damit auf 199 ppi.

\IfFileExists{../grafiken/stimmung-dorf.jpg}
  {\dsaStimmungsbild{\dsagrafik{stimmung-dorf}}}
  {\dsaStimmungsbild{\dsagrafik{ruecken-neutral}}}

%%% ---------------------------------------------------------------- 1
%%% Vier gleiche Spalten, der einfachste Fall.

\begin{dsaReihe}

\dsaKolumne{1}{%
  \begin{dsaEinlegertabelle}{Beispiele für Eigenschaftsmodifikatoren}{RW 19}{Xl}
    \dsaKopfzeile Modifikator & \dsaKopfschrift Bewertung \\
    \dsaMod{+}{6}  & extrem leichte Probe \\
    \dsaMod{+}{4}  & sehr leichte Probe \\
    \dsaMod{+}{2}  & leichte Probe \\
    \dsaMod{+}{}\slash\dsaMod{-}{ 0} & anspruchsvolle Probe \\
    \dsaMod{-}{2}  & schwere Probe \\
    \dsaMod{-}{4}  & sehr schwere Probe \\
    \dsaMod{-}{6}  & extrem schwere Probe \\
  \end{dsaEinlegertabelle}

  \dsaLuft{2}

  \begin{dsaEinlegertabelle}{Qualitätsstufen}{RW 23}{Xl}
    \dsaKopfzeile Fertigkeitspunkte & \dsaKopfschrift Qualitätsstufe \\
    0--3   & 1 \\
    4--6   & 2 \\
    7--9   & 3 \\
    10--12 & 4 \\
    13--15 & 5 \\
    16+    & 6 \\
  \end{dsaEinlegertabelle}
}

\dsaKolumne{1}{%
  \begin{dsaEinlegertabelle}{Ergebnisse bei Sammelproben}{RW 26}{lX}
    \dsaKopfzeile QS & \dsaKopfschrift Erfolg \\
    6  & Teilerfolg \\
    10 & Aufgabe erfüllt \\
  \end{dsaEinlegertabelle}

  \dsaLuft{2}

  \begin{dsaEinlegertabelle}{Anzahl der erlaubten Proben}{RW 26}{lX}
    \dsaKopfzeile Anzahl & \dsaKopfschrift Herausforderung \\
    5  & schwer \\
    7  & durchschnittlich \\
    10 & leicht \\
  \end{dsaEinlegertabelle}

  \dsaLuft{2}

  % Eine Tabelle mit mehreren Abschnitten: \dsaZwischenzeile setzt mitten
  % im Satz eine zweite Kopfzeile.
  \begin{dsaEinlegertabelle}{Modifikatoren für die Zeremonieprobe}{RW 312}{Xr}
    \dsaKopfzeile Ort\rlap{*} & \dsaKopfschrift Modifikator \\
    Heiligtum der Gottheit               & \dsaMod{+}{2} \\
    Tempel der Gottheit (geweihter Boden) & \dsaMod{+}{1} \\
    Unheiligtum eines Erzdämonen          & \dsaMod{-}{3} \\
    \dsaZwischenzeile Zeit\rlap{*} & \dsaKopfschrift Modifikator \\
    Monat des eigenen Gottes  & \dsaMod{+}{1} \\
    Feiertag des eigenen Gottes & \dsaMod{+}{2} \\
    Namenlose Tage            & \dsaMod{-}{5} \\
  \end{dsaEinlegertabelle}

  \dsaAnmerkung{*Nur einer der Modifikatoren kann gelten}
}

\dsaKolumne{1}{%
  % Eine Tabelle mit umbrechenden Zellen: T{} statt X, wenn die Breite fest
  % sein soll, X, wenn sie sich den Rest teilen.
  \begin{dsaEinlegertabelle}{Regeneration}{RW 339}{XL{22mm}}
    \dsaKopfzeile Situation & \dsaKopfschrift Regeneration \\
    \textbf{Grundregeneration für 6 Stunden Schlaf} & 1W6 LeP\slash AsP\slash KaP \\
    \textbf{Schlechter Lagerplatz, misslungene Probe auf Wildnisleben
      (Lagersuche)} & \dsaMod{-}{1} LeP\slash AsP\slash KaP \\
    \textbf{Unterbrechung der Nachtruhe (z.\,B. Wache halten\slash Ruhestörung)}
      & \dsaMod{-}{1} LeP\slash AsP\slash KaP \\
    \textbf{Held ist erkrankt\slash vergiftet} & keine LeP-Regeneration \\
    \textbf{Gute Unterkunft (Einzelzimmer Reisegasthof)}
      & \dsaMod{+}{1} LeP\slash AsP\slash KaP \\
  \end{dsaEinlegertabelle}
}

\dsaKolumne{1}{%
  \begin{dsaEinlegertabelle}{Routinevoraussetzungen}{RW 185}{L{24mm}r}
    \dsaKopfzeile[2] Modifikator-Maximum & \dsaKopfschrift Fertigkeitswert-Minimum \\
    \dsaMod{+}{3} und höher & 1 \\
    \dsaMod{+}{2}  & 4 \\
    \dsaMod{+}{1}  & 7 \\
    \dsaMod{+}{}\slash\dsaMod{-}{0} & 10 \\
    \dsaMod{-}{1}  & 13 \\
    \dsaMod{-}{2}  & 16 \\
    \dsaMod{-}{3}  & 19 \\
  \end{dsaEinlegertabelle}

  \dsaAnmerkung{*Alle beteiligten Eigenschaften mindestens 13}

  \dsaLuft{2}

  % Ein Pergamentkasten: derselbe Block, aber auf der Pergamentflaeche mit
  % gezacktem Rand statt auf blankem Grund. So steht er auch im Original auf
  % Seite 4. Innen gilt der normale Einlegersatz -- Blocktitel mit Marke,
  % Aufzaehlung, derselbe Schriftgrad.
  \begin{dsaPergament}{1}{9}
    \dsaBlocktitel[RW 29]{Einsatz von Schips}
    \begin{dsaliste}
      \item \textbf{Erster:} Der Held kann in einer Kampfrunde zuerst
        handeln.
      \item \textbf{Neuer Wurf:} Beliebig viele Würfel einer Probe dürfen
        erneut gewürfelt werden.
      \item \textbf{Qualität verbessern:} Die QS einer gelungenen Probe
        steigt um 1.
    \end{dsaliste}
  \end{dsaPergament}
}

\end{dsaReihe}

\newpage

%%% ---------------------------------------------------------------- 2
%%% Verschachtelte Reihen: links eine Spalte ueber die volle Hoehe, rechts
%%% eine breite Tabelle und darunter drei schmale. Das ist der Aufbau von
%%% Seite 4 des Originals.

\begin{dsaReihe}

\dsaKolumne{1}{%
  \begin{dsaEinlegertabelle}{Beispiele für Fertigkeitsmodifikatoren}{RW 23}{Xl}
    \dsaKopfzeile Modifikator & \dsaKopfschrift Bewertung \\
    \dsaMod{+}{5} & extrem leichte Probe \\
    \dsaMod{+}{3} & sehr leichte Probe \\
    \dsaMod{+}{1} & leichte Probe \\
    \dsaMod{+}{}\slash\dsaMod{-}{ 0} & anspruchsvolle Probe \\
    \dsaMod{-}{1} & schwierige Probe \\
    \dsaMod{-}{3} & sehr schwierige Probe \\
    \dsaMod{-}{5} & extrem schwierige Probe \\
  \end{dsaEinlegertabelle}
}

\dsaKolumne{3}{%
  % Eine Tabelle ueber drei Spalten. Sechs Wertespalten teilen sich die
  % Breite; die erste bleibt fest, damit die Zustandsnamen nicht umbrechen.
  \begin{dsaEinlegertabelle}{Zustände}{RW 31}{L{22mm}lXXXX}
    \dsaKopfzeile Zustand & \dsaKopfschrift Abbau
      & \dsaKopfschrift Stufe I & \dsaKopfschrift Stufe II
      & \dsaKopfschrift Stufe III & \dsaKopfschrift Stufe IV \\
    \textbf{Belastung} & -- & einige Proben auf Talente \dsaMinus 1; AT, VW, INI, GS \dsaMinus 1
      & einige Proben auf Talente \dsaMinus 2; AT, VW, INI, GS \dsaMinus 2
      & einige Proben auf Talente \dsaMinus 3; AT, VW, INI, GS \dsaMinus 3
      & \textit{Handlungsunfähig} \\
    \textbf{Betäubung} & 3 Std. & alle Proben \dsaMinus 1 & alle Proben \dsaMinus 2
      & alle Proben \dsaMinus 3 & \textit{Handlungsunfähig} \\
    \textbf{Furcht} & 5 Min. & alle Proben \dsaMinus 1 & alle Proben \dsaMinus 2
      & alle Proben \dsaMinus 3 & \textit{Handlungsunfähig} \\
    \textbf{Paralyse} & 30 Min. & Proben, die Bewegung erfordern, \dsaMinus 1, GS 75\,\%
      & Proben, die Bewegung erfordern, \dsaMinus 2, GS 50\,\%
      & Proben, die Bewegung erfordern, \dsaMinus 3, GS 25\,\%
      & \textit{Bewegungsunfähig} \\
    \textbf{Schmerz} & 4 Std. & alle Proben \dsaMinus 1, GS \dsaMinus 1
      & alle Proben \dsaMinus 2, GS \dsaMinus 2
      & alle Proben \dsaMinus 3, GS \dsaMinus 3
      & \textit{Handlungsunfähig}, ansonsten alle Proben \dsaMinus 4 \\
    \textbf{Verwirrung} & 1 Std. & alle Proben \dsaMinus 1 & alle Proben \dsaMinus 2
      & alle Proben \dsaMinus 3, komplexe Aktivitäten unmöglich
      & \textit{Handlungsunfähig} \\
  \end{dsaEinlegertabelle}

  \dsaLuft{2}

  % Verschachtelt: drei schmale Felder unter der breiten Tabelle.
  \begin{dsaReihe}
    \dsaKolumne{1}{%
      \begin{dsaEinlegertabelle}{Modifikatoren im Nahkampf}{RW 238}{Xl}
        \dsaKopfzeile Art & \dsaKopfschrift Modifikatoren \\
        \textbf{Kurze Waffen}   & \dsaMod{+}{}\slash\dsaMod{-}{0} auf AT \\
        \textbf{Mittlere Waffen} & \dsaMod{-}{4} auf AT \\
        \textbf{Lange Waffen}   & \dsaMod{-}{8} auf AT \\
        \textbf{Kleine Schilde} & \dsaMod{-}{2} auf AT \\
        \textbf{Große Schilde}  & \dsaMod{-}{6} auf AT \\
      \end{dsaEinlegertabelle}
    }
    \dsaKolumne{1}{%
      \begin{dsaEinlegertabelle}{Reichweite}{RW 241}{Xl}
        \dsaKopfzeile Entfernung & \dsaKopfschrift Modifikator \\
        \textbf{Nah}   & \dsaMod{+}{2} auf FK (\dsaMod{+}{1} TP) \\
        \textbf{Mittel} & \dsaMod{+}{}\slash\dsaMod{-}{0} auf FK \\
        \textbf{Weit}  & \dsaMod{-}{2} auf FK (\dsaMinus 1 TP) \\
      \end{dsaEinlegertabelle}
    }
    \dsaKolumne{1}{%
      \begin{dsaEinlegertabelle}{Größenkategorien}{RW 239}{L{16mm}Xl}
        \dsaKopfzeile Größe & \dsaKopfschrift Beispiele & \dsaKopfschrift Mod. \\
        \textbf{Winzig} & Ratte, Kröte, Spatz & \dsaMod{-}{4} AT \\
        \textbf{Klein}  & Rehkitz, Schaf, Ziege & \dsaMod{+}{}\slash\dsaMod{-}{0} AT \\
        \textbf{Mittel} & Mensch, Zwerg, Esel & \dsaMod{+}{}\slash\dsaMod{-}{0} AT \\
        \textbf{Groß}   & Oger, Troll, Rind & nur Schildparade \\
        \textbf{Riesig} & Drache, Riese, Elefant & nur Ausweichen \\
      \end{dsaEinlegertabelle}
    }
  \end{dsaReihe}
}

\end{dsaReihe}

\newpage

%%% ---------------------------------------------------------------- 3
%%% Ein Block ueber alle vier Spalten, darunter zwei zu je zweien und eine
%%% Ueberschrift ueber alles. Das ist der Aufbau der unteren Haelfte von
%%% Seite 6 des Originals, um die volle Breite ergaenzt.

%%% Eine Tabelle ueber die ganze Blattbreite: \dsaKolumne{4}. Sechs
%%% Wertespalten haben hier 275 mm zur Verfuegung statt 65.
\begin{dsaReihe}
  \dsaKolumne{4}{%
    \begin{dsaEinlegertabelle}{Zustände im Überblick}{RW 31}{L{26mm}lXXXX}
      \dsaKopfzeile Zustand & \dsaKopfschrift Abbau
        & \dsaKopfschrift Stufe I & \dsaKopfschrift Stufe II
        & \dsaKopfschrift Stufe III & \dsaKopfschrift Stufe IV \\
      \textbf{Berauscht} & 2 Std. & Proben auf Zechen \dsaMinus 1
        & Proben auf Zechen \dsaMinus 2 & Proben auf Zechen \dsaMinus 3
        & Held erleidet eine Stufe \textit{Betäubung} \\
      \textbf{Entrückung} & 1 Std. & einige Proben \dsaMinus 1
        & wohlgefällige Talente \dsaMod{+}{1}, einige Proben \dsaMinus 2
        & wohlgefällige Talente \dsaMod{+}{2}, einige Proben \dsaMinus 3
        & wohlgefällige Talente \dsaMod{+}{3}, einige Proben \dsaMinus 4 \\
      \textbf{Trance} & 24 Std. & keine AsP-Regeneration
        & keine AsP-Regeneration, alle Proben außer auf wohlgefällige
          Talente und Liturgien \dsaMinus 2
        & keine AsP-Regeneration, alle Proben \dsaMinus 3
        & \textit{Handlungsunfähig}, keine AsP-Regeneration \\
    \end{dsaEinlegertabelle}
  }
\end{dsaReihe}

\dsaLuft{3}

\begin{dsaReihe}
  \dsaKolumne{2}{%
    \begin{dsaEinlegertabelle}{Modifikatoren für die Ritualprobe}{RW 260}{Xr}
      \dsaKopfzeile Umstand & \dsaKopfschrift Modifikator \\
      Heiligtum oder magischer Ort & \dsaMod{+}{1} \\
      falscher Ritualplatz         & \dsaMod{-}{3} \\
      passende Konstellation (Mondstand meist Voll- oder Neumond)
        & \dsaMod{+}{1} \\
      unpassende Konstellation     & \dsaMod{-}{1} \\
      passende Kleidung            & \dsaMod{+}{1} \\
      besonders geeignete Gerätschaften & \dsaMod{+}{1} \\
    \end{dsaEinlegertabelle}
  }
  \dsaKolumne{2}{%
    \begin{dsaEinlegertabelle}{Magische Analyse}{RW 269}{lX}
      \dsaKopfzeile QS & \dsaKopfschrift Artefakt \slash{} Wesen \slash{} Zauberwirkung \\
      QS 1 & Eingrenzung der FP; mehr oder weniger als 10 FP; Bestimmung des
        Wesens: Dämon, Elementar, Geist \\
      QS 2 & Genauere Eingrenzung der FP; bis auf 3 FP genau \\
      QS 3 & Art des Artefakts (Zauberspeicher-Artefakt, Selbst aufladendes
        Artefakt, Dauerhaft wirkendes Artefakt) \\
      QS 4 & Tradition des Artefakterschaffers \slash{} Tradition des
        Beschwörers \\
    \end{dsaEinlegertabelle}
  }
\end{dsaReihe}

\dsaLuft{3}

\dsaUeberschrift{Modifikationen bei Zaubersprüchen und Liturgien
  sowie Ritualen und Zeremonien}

\dsaLuft{2}

\begin{dsaReihe}
  \dsaKolumne{2}{%
    \dsaBlocktitel[RW 255/309]{Modifikationen im Überblick}
    \begin{dsaliste}
      \item \textbf{Erzwingen:} Kosten um eine Stufe erhöht, Probe \dsaMod{+}{1}
      \item \textbf{Kosten senken:} Kosten um einen Schritt gesenkt,
        Probe \dsaMinus 1
      \item \textbf{Reichweite erhöhen:} Reichweite um eine Stufe erhöht,
        Probe \dsaMinus 1
      \item \textbf{Zauber-\slash Liturgiedauer erhöhen:} Dauer um eine Stufe
        erhöht, Probe \dsaMod{+}{1}
      \item \textbf{Geste oder Formel weglassen:} Jeweils Probe \dsaMinus 2
    \end{dsaliste}
  }
  \dsaKolumne{2}{%
    \begin{dsaEinlegertabelle}{Modifikationstabelle Zaubersprüche
      und Liturgien}{RW 255/309}{lXXXXX}
      \dsaKopfzeile Dauer & \dsaKopfschrift 1 Aktion
        & \dsaKopfschrift 2 Akt. & \dsaKopfschrift 4 Akt.
        & \dsaKopfschrift 8 Akt. & \dsaKopfschrift 16 Akt. \\
      \textbf{Reichweite} & Berühren & 4 Schritt & 8 Schritt & 16 Schritt
        & 32 Schritt \\
      \textbf{Kosten} & 1 AsP\slash KaP & 2 AsP\slash KaP & 4 AsP\slash KaP
        & 8 AsP\slash KaP & 16 AsP\slash KaP \\
    \end{dsaEinlegertabelle}

    \dsaLuft{2}

    \begin{dsaEinlegertabelle}{Modifikationstabelle Rituale
      und Zeremonien}{RW 260/313}{lXXXXX}
      \dsaKopfzeile Dauer & \dsaKopfschrift 5 Min.
        & \dsaKopfschrift 30 Min. & \dsaKopfschrift 2 Std.
        & \dsaKopfschrift 8 Std. & \dsaKopfschrift 16 Std. \\
      \textbf{Reichweite} & Berühren & 4 Schritt & 8 Schritt & 16 Schritt
        & 32 Schritt \\
      \textbf{Kosten} & 8 AsP\slash KaP & 16 AsP\slash KaP & 32 AsP\slash KaP
        & 64 AsP\slash KaP & 128 AsP\slash KaP \\
    \end{dsaEinlegertabelle}
  }
\end{dsaReihe}


\dsaLuft{3}


%%% ---------------------------------------------------------------- 4
%%% Der Rest der Elemente: Blocktitel ohne Marke, Marke allein, Tabelle
%%% ohne Titel, zentrierte Spalte, selbst gerechnete Breite.

\begin{dsaReihe}
  \dsaKolumne{1}{%
    % Blocktitel ohne Quellenmarke.
    \dsaBlocktitel{Ohne Quellenangabe}
    Ein Block muss keine Marke tragen. Der Titel steht dann allein an der
    Spaltenkante.
  }
  \dsaKolumne{1}{%
    % Tabelle ohne Titelzeile: erstes Argument leer lassen. Nuetzlich fuer
    % den zweiten Teil einer Tabelle, die in der Spalte daneben weitergeht.
    % Eine X-Spalte muss dabei sein, sonst erreicht die Tabelle die
    % Kolumnenbreite nicht.
    \begin{dsaEinlegertabelle}{}{}{Xcr}
      \dsaKopfzeile links & \dsaKopfschrift Mitte & \dsaKopfschrift rechts \\
      a & b & c \\
      d & e & f \\
    \end{dsaEinlegertabelle}
  }
  \dsaKolumne{2}{%
    % Die Quellenmarke allein, mitten im Text.
    \dsaBlocktitel{Marke und eigene Breite}
    Die Marke \dsaQuelle{RW 12} steht auch allein im Text.

    \dsaLuft{1}
    % \dsaKolumnenbreite{n} ist die Breite von n Spalten samt der Abstaende
    % dazwischen. Es ist ein \dimexpr und taugt nicht unmittelbar als
    % Argument von minipage -- erst in eine Laenge, dann einsetzen.
    %
    % \verb geht hier nicht: der Inhalt einer \dsaKolumne ist ein
    % Makroargument, und \verb liest seinen Text zeichenweise aus dem
    % Eingabestrom. Der Lauf meldet "\verb illegal in argument".
    \setlength{\dsahilfe}{\dsaKolumnenbreite{1}}%
    \noindent\begin{minipage}{\dsahilfe}%
      Diese Box ist genau eine Spalte breit, gerechnet mit
      \texttt{\textbackslash dsaKolumnenbreite\{1\}}.%
    \end{minipage}
  }
\end{dsaReihe}

\end{document}
